\chapter{Variabili aleatorie continue: densità e distribuzione}

Finora i risultati possibili si potevano mettere in fila e contare: le facce di un dado, le mani di carte, i pacchetti
in coda. Molte grandezze fisiche però variano con continuità: una tensione, un tempo di attesa, la posizione di un
oggetto. Una tensione può valere $0{,}31$ volt, o $0{,}312$, o $0{,}3127$\ldots{} e fra due valori qualsiasi ce ne sono
infiniti altri. In questo capitolo scopriamo che in questi casi la probabilità di un singolo valore è sempre
\textbf{zero}, e che serve un oggetto nuovo per descrivere la variabile: la \textbf{densità di probabilità}. Per
lavorarci useremo lo strumento matematico che fa per le grandezze continue quello che la sommatoria fa per quelle
discrete: l'\textbf{integrale}, spiegato da zero.

% ============================================================
\section{Spazi campione continui}
% ============================================================

Rimuoviamo ora l'ipotesi che lo spazio campione sia discreto. Supponiamo d'ora in poi che $\Omega \subseteq \R$ sia un
sottoinsieme \textbf{continuo} dei numeri reali (per esempio un intervallo). $\Omega$ può essere esso stesso lo spazio
delle misure osservabili, oppure il dominio di una variabile aleatoria
\[ X : \omega \in \Omega \;\longmapsto\; X(\omega) \in \mathcal{X} , \]
e ora l'alfabeto $\mathcal{X}$ non potrà più essere un insieme discreto. Spesso accade che $X(\omega) = \omega$ e
$\mathcal{X} = \Omega$: il risultato dell'esperimento è già il numero che interessa.

\begin{esempio}{Il rumore termico}
La tensione misurata a vuoto ai capi di un resistore non è mai esattamente zero: gli elettroni del materiale si
agitano a causa della temperatura (è il \emph{rumore termico}), e la tensione oscilla attorno allo zero in modo
imprevedibile.
\begin{enumerate}[label=\alph*)]
  \item Si misura $n$ volte questa tensione. Il risultato è già un numero reale: $X(\omega) = \omega = x \in \R$, e le
        misure sono $n$ numeri $\{x_i\}_{i=1}^{n}$.
  \item Si misura invece la potenza che la tensione trasferisce a un carico di resistenza $R$. Lo spazio campione è
        ancora lo stesso, ma la variabile aleatoria ora è $X(\omega) = \omega^2/R$: una funzione del risultato.
\end{enumerate}
\end{esempio}

\subsection{Perché un singolo valore ha probabilità zero}

In entrambi i casi gli eventi elementari sono del tipo $\{X(\omega_i) = x_i\}$. Si potrebbe essere tentati di
definire, come nel caso discreto,
\[ P(X = x_i) = \lim_{n \to \infty} \frac{n_{X = x_i}}{n} , \]
dove $n_{X = x_i}$ è il numero di volte in cui si è osservato esattamente il valore $x_i$. La definizione non è
sbagliata, ma è inutile. Due realizzazioni distinte $X(\omega_i)$ e $X(\omega_j)$ di una variabile aleatoria reale non
si possono mai misurare con esattezza: servirebbe uno strumento di precisione infinita, e anche così l'evento
$X(\omega_i) = X(\omega_j)$, che i due valori coincidano fino all'ultima delle infinite cifre decimali, sarebbe
impossibile. Quello che possiamo dire è se la misura $X(\omega_j)$ cade o meno \textbf{in un intorno} di
$X(\omega_i)$, cioè abbastanza vicino. Quindi, se $X$ è una variabile aleatoria continua, gli eventi elementari
$\{X = x\}$ hanno (a meno di casi degeneri) probabilità nulla.

C'è anche una ragione matematica: i valori possibili sono troppi perché ognuno possa avere una probabilità positiva.
Prima di vederla, un richiamo che chiarisce quando una somma di infiniti termini positivi può dare un numero finito.

\begin{richiamo}{la serie geometrica}
Per un numero $q$ con $\abs{q} < 1$ (per esempio $q = \frac12$, oppure $q = -0{,}3$):
\[ \sum_{k=0}^{\infty} q^k = 1 + q + q^2 + q^3 + \cdots = \frac{1}{1 - q} . \]
\textbf{Dimostrazione.} Sia $S_N = 1 + q + q^2 + \cdots + q^N$ la somma parziale. Moltiplicandola per $q$ si ottiene
$q\,S_N = q + q^2 + \cdots + q^{N+1}$: gli stessi termini, spostati di un posto. Sottraendo, quasi tutto si cancella
(è una somma telescopica):
\[ S_N - q\,S_N = 1 - q^{N+1} \quad\Longrightarrow\quad S_N = \frac{1 - q^{N+1}}{1 - q} . \]
Se $\abs{q} < 1$, le potenze $q^{N+1}$ diventano piccolissime al crescere di $N$ (ogni volta si moltiplica per un
numero più piccolo di 1 in valore assoluto), cioè $q^{N+1} \to 0$, e quindi $S_N \to \frac{1}{1-q}$.

Con $q = \frac12$ si ritrova $1 + \frac12 + \frac14 + \cdots = 2$ (la serie del Capitolo~4 partiva da $\frac12$ e valeva
1). Una somma di \emph{infiniti} termini positivi, quindi, può essere finita, purché i termini diventino piccoli abbastanza
in fretta. Questo però succede per un'infinità \textbf{numerabile} di termini, quelli che si possono mettere in fila.
\end{richiamo}

\teorema[atomi]{Pochi valori possono avere probabilità positiva}{
  Per qualsiasi variabile aleatoria, l'insieme dei valori $x$ con $P(X = x) > 0$ è al più numerabile. In particolare,
  se l'alfabeto è un intervallo di numeri reali, non tutti i suoi valori possono avere probabilità positiva: quasi
  tutti devono avere probabilità zero.
}

\begin{dimostrazione}
\passo{Pochi valori «grossi»} Fissato un intero $k \geq 1$, i valori con $P(X = x) \geq \frac1k$ sono al massimo $k$:
se fossero $k + 1$, gli eventi $\{X = x\}$ corrispondenti, che sono incompatibili, avrebbero probabilità totale
$\geq \frac{k+1}{k} > 1$. Per esempio, al massimo 2 valori possono avere probabilità $\geq \frac12$, al massimo 10
valori probabilità $\geq \frac1{10}$.
\passo{Li si mette in fila} Un valore con probabilità positiva ha probabilità $\geq \frac1k$ per qualche $k$ (basta
prendere $k$ abbastanza grande). Quindi i valori con probabilità positiva si possono elencare: prima quelli (al più
uno) con probabilità $\geq 1$, poi quelli (al più due) con probabilità $\geq \frac12$, poi quelli $\geq \frac13$, e così
via. L'elenco può essere infinito, ma è un elenco: l'insieme è al più numerabile.
\passo{L'intervallo è troppo grande} I numeri reali di un intervallo non sono numerabili (Capitolo~1): non si possono
mettere tutti in un elenco. Quindi non tutti possono avere probabilità positiva.
\end{dimostrazione}

In una variabile continua «onesta» nessun valore è privilegiato rispetto ai suoi vicini, e allora tutti hanno
probabilità zero. Ma attenzione: probabilità zero \textbf{non} vuol dire impossibile. Ogni volta che si misura una
tensione si osserva \emph{qualche} valore, e quel valore aveva probabilità zero di uscire.

\subsection{La frequenza di un intervallo}

Se la probabilità dei singoli valori è nulla, ha senso solo chiedersi con che frequenza la variabile cade in un
\textbf{intervallo}. Supponiamo di compiere $n$ esperimenti e di disporre delle osservazioni $X(\omega_i)$. Fissato un
valore $x$, contiamo quante volte $X$ è caduta in un intorno di $x$ di ampiezza $\Delta x$, cioè fra
$x - \frac{\Delta x}{2}$ e $x + \frac{\Delta x}{2}$:
\[ f_n(x; \Delta x) = \frac{n_{\left\{x - \frac{\Delta x}{2} \,\leq\, X \,\leq\, x + \frac{\Delta x}{2}\right\}}}{n} . \]
Possiamo allora definire nel modo consueto la probabilità dell'evento «$X$ cade nell'intervallo»:
\[
  P\left( \left\{ \omega \in \Omega : x - \tfrac{\Delta x}{2} \leq X \leq x + \tfrac{\Delta x}{2} \right\} \right)
  = P\left( X \in \left[ x - \tfrac{\Delta x}{2},\ x + \tfrac{\Delta x}{2} \right] \right)
  = P_X(x; \Delta x) = \lim_{n \to \infty} f_n(x; \Delta x) .
\]

% ============================================================
\section{La densità di probabilità}
% ============================================================

La probabilità di un intervallo dipende dalla sua ampiezza: un intervallo largo il doppio raccoglie, più o meno, il
doppio della probabilità. Per avere una grandezza che descriva la variabile \emph{in un punto} si divide la
probabilità per l'ampiezza, ottenendo una probabilità \emph{per unità di lunghezza}, e poi si stringe l'intervallo.

\dfn{Densità di probabilità (PDF)}{
  La \term{funzione densità di probabilità} (\emph{probability density function}, \term{PDF}) della variabile
  aleatoria continua $X$ è
  \[ f_X(x) = \lim_{\Delta x \to 0} \frac{P\left( x - \frac{\Delta x}{2} \leq X \leq x + \frac{\Delta x}{2}
     \right)}{\Delta x} = \lim_{\Delta x \to 0} \frac{P_X(x; \Delta x)}{\Delta x} . \]
}

Il nome viene dalla fisica. La densità di una sbarra di metallo è la sua massa per unità di lunghezza: una sbarra
più densa in un punto pesa di più in quel tratto. Allo stesso modo $f_X(x)$ dice quanta probabilità c'è «per ogni
centimetro» attorno a $x$. Nel caso discreto la probabilità era concentrata in mucchietti (i bastoncini della pmf);
ora è spalmata con continuità, come uno strato di sabbia più spesso in certi punti e più sottile in altri.

\begin{figure}[H]
\centering
\begin{tikzpicture}
\begin{groupplot}[group style={group size=3 by 1, horizontal sep=0.6cm},
  width=5.6cm, height=4.6cm, xmin=-4, xmax=4, ymin=0, ymax=0.45, xtick={-3,0,3}, ytick={0,0.2,0.4},
  yticklabels={$0$,$0{,}2$,$0{,}4$}, axis lines=left, tick label style={font=\scriptsize},
  title style={font=\small}]
\nextgroupplot[title={$n = 100$, $\Delta x = 1$}]
  \addplot[ybar interval, fill=mspiA!40, draw=mspiA!85!black] coordinates {(-4.00,0.0100) (-3.00,0.0100) (-2.00,0.1000) (-1.00,0.3200) (0.00,0.3900) (1.00,0.1500) (2.00,0.0200) (3.00,0.0000) (4.00,0.0000)};
  \addplot[mspiB, line width=1pt, domain=-4:4, samples=80] {exp(-x^2/2)/sqrt(2*pi)};
\nextgroupplot[title={$n = 1000$, $\Delta x = 0{,}5$}, yticklabels={}]
  \addplot[ybar interval, fill=mspiA!40, draw=mspiA!85!black] coordinates {(-4.00,0.0000) (-3.50,0.0020) (-3.00,0.0100) (-2.50,0.0340) (-2.00,0.1040) (-1.50,0.2140) (-1.00,0.2960) (-0.50,0.3660) (0.00,0.3600) (0.50,0.3300) (1.00,0.1780) (1.50,0.0700) (2.00,0.0260) (2.50,0.0100) (3.00,0.0000) (3.50,0.0000) (4.00,0.0000)};
  \addplot[mspiB, line width=1pt, domain=-4:4, samples=80] {exp(-x^2/2)/sqrt(2*pi)};
\nextgroupplot[title={$n = 100\,000$, $\Delta x = 0{,}1$}, yticklabels={}]
  \addplot[ybar interval, fill=mspiA!40, draw=mspiA!70, line width=0.2pt] coordinates {(-4.00,0.0002) (-3.90,0.0003) (-3.80,0.0005) (-3.70,0.0004) (-3.60,0.0003) (-3.50,0.0007) (-3.40,0.0017) (-3.30,0.0023) (-3.20,0.0033) (-3.10,0.0038) (-3.00,0.0056) (-2.90,0.0074) (-2.80,0.0092) (-2.70,0.0125) (-2.60,0.0141) (-2.50,0.0215) (-2.40,0.0255) (-2.30,0.0329) (-2.20,0.0415) (-2.10,0.0505) (-2.00,0.0574) (-1.90,0.0752) (-1.80,0.0896) (-1.70,0.0994) (-1.60,0.1217) (-1.50,0.1374) (-1.40,0.1572) (-1.30,0.1807) (-1.20,0.2070) (-1.10,0.2291) (-1.00,0.2430) (-0.90,0.2846) (-0.80,0.3012) (-0.70,0.3170) (-0.60,0.3398) (-0.50,0.3511) (-0.40,0.3647) (-0.30,0.3858) (-0.20,0.3889) (-0.10,0.3993) (0.00,0.3994) (0.10,0.3883) (0.20,0.3926) (0.30,0.3772) (0.40,0.3552) (0.50,0.3496) (0.60,0.3283) (0.70,0.3052) (0.80,0.2774) (0.90,0.2564) (1.00,0.2366) (1.10,0.2070) (1.20,0.1946) (1.30,0.1610) (1.40,0.1384) (1.50,0.1229) (1.60,0.0994) (1.70,0.0878) (1.80,0.0720) (1.90,0.0612) (2.00,0.0521) (2.10,0.0389) (2.20,0.0312) (2.30,0.0227) (2.40,0.0189) (2.50,0.0143) (2.60,0.0121) (2.70,0.0098) (2.80,0.0058) (2.90,0.0057) (3.00,0.0039) (3.10,0.0029) (3.20,0.0015) (3.30,0.0016) (3.40,0.0012) (3.50,0.0011) (3.60,0.0004) (3.70,0.0003) (3.80,0.0003) (3.90,0.0002) (4.00,0.0002)};
  \addplot[mspiB, line width=1pt, domain=-4:4, samples=80] {exp(-x^2/2)/sqrt(2*pi)};
\end{groupplot}
\end{tikzpicture}
\caption{Istogrammi di misure di rumore (simulate): ogni barra è alta $f_n(x; \Delta x)/\Delta x$, la frequenza
dell'intervallo divisa per la sua ampiezza. Con più misure e intervalli più stretti l'istogramma si avvicina a una curva
liscia (in arancione): la densità di probabilità. Questa curva è la gaussiana, che studieremo nel Capitolo~15.}
\label{fig:istogrammi}
\end{figure}

\subsection{Dalla densità alla probabilità: l'integrale}

La definizione dice che, per un intervallo piccolo, la probabilità è circa la densità per l'ampiezza:
\[ P\left( x - \tfrac{\Delta x}{2} \leq X \leq x + \tfrac{\Delta x}{2} \right) \simeq f_X(x)\, \Delta x . \]
Per un intervallo grande $[a_1, a_2]$ lo si spezza in tanti intervalli piccoli, si calcola la probabilità di ciascuno
così e si somma: in un grafico, è la somma delle aree di tanti rettangoli sottili sotto la curva $f_X$. Quando i
rettangoli diventano infinitamente sottili, la somma diventa l'\textbf{area sotto la curva}, che si chiama integrale.

\begin{richiamo}{l'integrale come area}
Sia $f(x) \geq 0$ una funzione definita sull'intervallo $[a, b]$. Si vuole calcolare l'area della regione compresa fra
il grafico di $f$ e l'asse $x$, fra $a$ e $b$. Si divide $[a, b]$ in $n$ pezzi uguali di larghezza
$\Delta = \frac{b-a}{n}$, con estremi $x_0 = a, x_1, \ldots, x_n = b$, e su ogni pezzo si costruisce un rettangolo alto
quanto la funzione in un punto del pezzo. La somma delle aree dei rettangoli,
\[ \sum_{i=0}^{n-1} f(x_i)\, \Delta , \]
si chiama \term{somma di Riemann}. Più i rettangoli sono sottili, meglio approssimano l'area:
\begin{center}
\begin{tikzpicture}
\begin{groupplot}[group style={group size=2 by 1, horizontal sep=1.2cm},
  width=6.4cm, height=4.4cm, xmin=0, xmax=5, ymin=0, ymax=2.4, xtick={0.5,4.5}, xticklabels={$a$,$b$},
  ytick=\empty, axis lines=left, tick label style={font=\small}, title style={font=\small}]
\nextgroupplot[title={4 rettangoli}]
  \addplot[ybar interval, fill=mspiric!25, draw=mspiric!80!black, samples=5, domain=0.5:4.5] {1.3+0.7*sin(deg(1.1*x))};
  \addplot[mspiA, line width=1.1pt, domain=0.3:4.7, samples=60] {1.3+0.7*sin(deg(1.1*x))};
\nextgroupplot[title={16 rettangoli}]
  \addplot[ybar interval, fill=mspiric!25, draw=mspiric!80!black, samples=17, domain=0.5:4.5] {1.3+0.7*sin(deg(1.1*x))};
  \addplot[mspiA, line width=1.1pt, domain=0.3:4.7, samples=60] {1.3+0.7*sin(deg(1.1*x))};
\end{groupplot}
\end{tikzpicture}
\end{center}
Il limite delle somme di Riemann per $n \to \infty$ (cioè per $\Delta \to 0$) è l'\term{integrale definito}:
\[ \int_a^b f(x) \dd x = \lim_{\Delta \to 0} \sum_{i} f(x_i)\, \Delta . \]
La notazione racconta la costruzione: il simbolo $\int$ è una «S» allungata (di somma), $f(x)$ è l'altezza di un
rettangolo e $\mathrm{d}x$ è la sua larghezza infinitamente piccola. La lettera $x$ è una variabile muta, come l'indice
di una sommatoria: $\int_a^b f(x) \dd x = \int_a^b f(t) \dd t$.

Proprietà che seguono subito dalla costruzione con i rettangoli:
\begin{itemize}
  \item \textbf{linearità}: $\int_a^b \left( c\,f(x) + d\,g(x) \right) \dd x = c \int_a^b f(x) \dd x + d \int_a^b g(x) \dd x$;
  \item \textbf{additività}: se $a < c < b$, $\int_a^b f = \int_a^c f + \int_c^b f$ (l'area si può tagliare in due);
  \item \textbf{positività}: se $f(x) \geq 0$ su $[a, b]$, allora $\int_a^b f(x) \dd x \geq 0$;
  \item un singolo punto non ha area: $\int_a^a f(x) \dd x = 0$.
\end{itemize}
Se la funzione è negativa in qualche tratto, le aree sotto l'asse $x$ contano con il segno meno.
\end{richiamo}

Con l'integrale, la probabilità di qualsiasi intervallo si ottiene dalla densità.

\begin{formula}[probabilità di un intervallo]
\[ P(a_1 \leq X \leq a_2) = \int_{a_1}^{a_2} f_X(t) \dd t \qquad\qquad
   P\left( x - \tfrac{\Delta x}{2} \leq X \leq x + \tfrac{\Delta x}{2} \right) = \int_{x - \frac{\Delta x}{2}}^{x +
   \frac{\Delta x}{2}} f_X(t) \dd t \]
\end{formula}

\begin{figure}[H]
\centering
\begin{tikzpicture}
\begin{axis}[quaderno, width=10cm, height=5cm, xmin=-3.6, xmax=3.6, ymin=0, ymax=0.48, xtick={-0.6,1.4},
  xticklabels={$a_1$,$a_2$}, ytick=\empty, xlabel=$x$, ylabel=$f_X(x)$, grid=none]
  \addplot[fill=evid, draw=none, domain=-0.6:1.4, samples=40] {exp(-x^2/2)/sqrt(2*pi)} \closedcycle;
  \addplot[mspiA, domain=-3.4:3.4, samples=80] {exp(-x^2/2)/sqrt(2*pi)};
  \node[font=\small, align=center] at (axis cs:0.4,0.16) {area $=$\\ $P(a_1 \leq X \leq a_2)$};
  \node[font=\small, mspiA!85!black, anchor=west] at (axis cs:1.9,0.2) {area totale $= 1$};
\end{axis}
\end{tikzpicture}
\caption{La probabilità che $X$ cada fra $a_1$ e $a_2$ è l'area sotto la densità, fra $a_1$ e $a_2$.}
\end{figure}

\subsection{Le proprietà della densità}

Per l'area «totale», da $-\infty$ a $+\infty$, serve un'estensione dell'integrale.

\begin{richiamo}{gli integrali impropri}
Un integrale su un intervallo infinito si definisce come limite:
\[ \int_a^{\infty} f(x) \dd x = \lim_{H \to \infty} \int_a^H f(x) \dd x , \qquad
   \int_{-\infty}^{\infty} f(x) \dd x = \lim_{H \to \infty} \int_{-H}^{0} f(x) \dd x + \lim_{K \to \infty} \int_0^K
   f(x) \dd x . \]
Se i limiti esistono e sono finiti, l'integrale si dice \term{convergente} (la funzione è \term{sommabile}). Una regione
infinitamente lunga può avere area finita, se la funzione si abbassa abbastanza in fretta: è la stessa idea della
serie geometrica. Per esempio (lo calcoleremo con il teorema del prossimo richiamo)
$\int_1^H \frac{1}{x^2} \dd x = 1 - \frac1H$, che tende a 1: l'area sotto $\frac{1}{x^2}$ da 1 a infinito vale 1.
Invece $\int_0^H 1 \dd x = H$ cresce senza limite: un rettangolo alto 1 e infinitamente lungo ha area infinita.

Nel secondo caso le due metà devono convergere \emph{ciascuna per conto suo}: non basta che si compensino. È un punto
che diventerà importante con la variabile di Cauchy, nel prossimo capitolo.
\end{richiamo}

\teorema[proprietapdf]{Proprietà della densità}{
  Ogni densità di probabilità soddisfa:
  \begin{enumerate}[label=\alph*)]
    \item $f_X(x) \geq 0$ per ogni $x \in \R$;
    \item $f_X$ è sommabile su $\R$ e il suo integrale vale 1:
          $\displaystyle \int_{-\infty}^{\infty} f_X(t) \dd t = P(X \in \R) = P(\Omega) = 1$.
  \end{enumerate}
  Inoltre $P(X = x) = 0$ per ogni $x$, quindi gli estremi di un intervallo non contano:
  $P(a_1 \leq X \leq a_2) = P(a_1 < X < a_2) = P(a_1 \leq X < a_2)$.
}

\begin{dimostrazione}
\passo{Non negatività} $f_X(x)$ è il limite di rapporti fra una probabilità ($\geq 0$) e un'ampiezza ($> 0$). In altre
parole: se la densità fosse negativa in un tratto, il suo integrale su quel tratto, cioè la probabilità di
quell'intervallo, sarebbe negativo.
\passo{Area totale} L'evento «$X$ è un numero reale qualsiasi» è certo, quindi ha probabilità 1, ed è l'area sotto
tutta la curva.
\passo{Un punto} $P(X = x) \leq P\left( x - \frac{\Delta x}{2} \leq X \leq x + \frac{\Delta x}{2} \right) \simeq
f_X(x)\, \Delta x$, che tende a zero quando $\Delta x \to 0$. Quindi aggiungere o togliere gli estremi a un intervallo
non cambia la sua probabilità.
\end{dimostrazione}

\warning{Una densità non è una probabilità}{
  $f_X(x)$ può benissimo essere \textbf{maggiore di 1}. Una variabile uniformemente distribuita fra 0 e $\frac12$ ha
  densità 2 su quell'intervallo: l'area è $2 \cdot \frac12 = 1$, come deve. Ad avere un significato di probabilità è
  l'area, non l'altezza: $f_X(x)\,\Delta x$, per $\Delta x$ piccolo. Ed è $P(X = x) = 0$, non $f_X(x)$.
}

% ============================================================
\section{La funzione di distribuzione cumulativa}
% ============================================================

La densità $f_X(x)$, $x \in \R$, è perfettamente adeguata a caratterizzare la variabile aleatoria $X$. I punti in cui
è diversa da zero formano il suo \term{supporto}, $\mathrm{supp}[f_X(x)] = \mathcal{X}$: sono i valori che la
variabile può assumere. E la probabilità di qualsiasi intervallo è l'integrale della densità. Tuttavia è molto usata
anche una caratterizzazione alternativa, che parla direttamente di probabilità.

\dfn{Funzione di distribuzione cumulativa (CDF)}{
  La \term{funzione di distribuzione cumulativa} (\emph{Cumulative Distribution Function}, \term{CDF}) di $X$ è la
  probabilità che $X$ non superi il valore $x$:
  \[ F_X(x) = P(-\infty < X \leq x) = \int_{-\infty}^{x} f_X(t) \dd t . \]
}

$F_X(x)$ è l'area sotto la densità \emph{a sinistra} di $x$: la probabilità «accumulata» da $-\infty$ fino a $x$.
Quando $x$ si sposta verso destra, l'area cresce. Il legame fra $F_X$ e $f_X$ è uno dei risultati più importanti
della matematica.

\begin{richiamo}{il teorema fondamentale del calcolo integrale}
Sia $f$ una funzione continua e sia $F(x) = \int_a^x f(t) \dd t$ la \term{funzione integrale}: l'area sotto $f$ da
$a$ fino a $x$, vista come funzione dell'estremo destro. Allora
\[ F'(x) = f(x) : \]
\textbf{la derivata dell'area è l'altezza della curva}. Il motivo si vede nel disegno: quando l'estremo passa da $x$ a
$x + h$, l'area cresce di una strisciolina larga $h$ e alta circa $f(x)$.
\begin{center}
\begin{tikzpicture}
\begin{axis}[quaderno, width=8cm, height=4.6cm, xmin=0, xmax=5, ymin=0, ymax=2.4, xtick={0.5,3,3.5},
  xticklabels={$a$,$x$,$x+h$}, ytick=\empty, grid=none]
  \addplot[fill=mspiric!20, draw=none, domain=0.5:3, samples=40] {1.3+0.7*sin(deg(1.1*x))} \closedcycle;
  \addplot[fill=evid, draw=none, domain=3:3.5, samples=10] {1.3+0.7*sin(deg(1.1*x))} \closedcycle;
  \addplot[mspiA, domain=0.3:4.7, samples=60] {1.3+0.7*sin(deg(1.1*x))};
  \node[font=\small] at (axis cs:1.75,0.55) {$F(x)$};
  \draw[-{Stealth}] (axis cs:4.2,1.9) node[right, font=\scriptsize, align=left] {striscia: area\\ $\simeq f(x)\,h$}
    -- (axis cs:3.3,0.9);
\end{axis}
\end{tikzpicture}
\end{center}
Quindi $F(x + h) - F(x) \simeq f(x)\, h$, il rapporto incrementale $\frac{F(x+h) - F(x)}{h}$ è circa $f(x)$, e diventa
esattamente $f(x)$ quando $h \to 0$ (la striscia diventa un rettangolo perfetto).

\textbf{Come si calcola un integrale.} Una funzione $G$ con $G'(x) = f(x)$ si chiama \term{primitiva} di $f$. Due
primitive della stessa funzione differiscono per una costante (la loro differenza ha derivata nulla, quindi è
piatta). Siccome anche $F$ è una primitiva, $F(x) = G(x) + c$, e da $F(a) = 0$ si ricava $c = -G(a)$. Quindi
\[ \int_a^b f(x) \dd x = G(b) - G(a) = \Big[ G(x) \Big]_a^b : \]
per calcolare un integrale basta trovare una primitiva e fare la differenza dei suoi valori agli estremi. Le primitive
si trovano leggendo al contrario la tabella delle derivate. Da $(x^{n+1})' = (n+1)\,x^n$:
\[ \int_a^b x^n \dd x = \left[ \frac{x^{n+1}}{n+1} \right]_a^b = \frac{b^{n+1} - a^{n+1}}{n+1} \qquad (n \neq -1) . \]
Per esempio $\int_0^1 x \dd x = \frac12$ (l'area del triangolo con base e altezza 1) e
$\int_1^H x^{-2} \dd x = \left[ -x^{-1} \right]_1^H = 1 - \frac1H$.
\end{richiamo}

Applicando il teorema fondamentale alla CDF si ottiene subito il legame con la densità.

\begin{formula}[CDF e PDF]
\[ F_X(x) = \int_{-\infty}^{x} f_X(t) \dd t \qquad \Longleftrightarrow \qquad f_X(x) = \frac{\dd}{\dd x} F_X(x) \]
\end{formula}

Talvolta si fa riferimento alla funzione complementare, la \term{CCDF} (\emph{Complementary Cumulative Distribution
Function}): la probabilità che $X$ superi $x$, cioè l'area a destra di $x$,
\[ \overline{F}_X(x) = P(X > x) = \int_x^{\infty} f_X(t) \dd t = 1 - F_X(x)
   \quad\Longrightarrow\quad f_X(x) = -\frac{\dd}{\dd x}\, \overline{F}_X(x) . \]

\teorema[proprietacdf]{Proprietà della CDF}{
  \begin{enumerate}[label=\alph*)]
    \item $F_X(x) \in [0, 1]$;
    \item $F_X(-\infty) = 0$ e $F_X(+\infty) = 1$ (intesi come limiti);
    \item $F_X$ è continua;
    \item $F_X$ è crescente (non decrescente): se $x_1 < x_2$ allora $F_X(x_1) \leq F_X(x_2)$;
    \item $P(a_1 \leq X \leq a_2) = \int_{a_1}^{a_2} f_X(t) \dd t = F_X(a_2) - F_X(a_1)$.
  \end{enumerate}
}

\begin{dimostrazione}
\passo{Fra 0 e 1} $F_X(x)$ è una probabilità.
\passo{Ai due estremi} Per $x \to -\infty$ l'area a sinistra di $x$ si riduce a nulla; per $x \to +\infty$ diventa
tutta l'area sotto la densità, che vale 1.
\passo{Continuità} Spostando $x$ di poco, l'area cambia di poco: di una striscia stretta, larga quanto lo spostamento.
(È la funzione integrale di una funzione sommabile.)
\passo{Crescenza} Se $x_1 < x_2$, l'area fino a $x_2$ è quella fino a $x_1$ più l'area fra $x_1$ e $x_2$, che non è
negativa perché l'integrando $f_X \geq 0$.
\passo{Un intervallo} Per l'additività dell'integrale, $\int_{-\infty}^{a_2} f_X = \int_{-\infty}^{a_1} f_X +
\int_{a_1}^{a_2} f_X$, cioè $F_X(a_2) = F_X(a_1) + P(a_1 \leq X \leq a_2)$.
\end{dimostrazione}

\begin{esempio}{La densità triangolare}
Sia $X$ con densità a forma di triangolo:
\[ f_X(x) = \begin{cases} 1 - \abs{x} & -1 \leq x \leq 1 \\ 0 & \text{altrove} \end{cases} \]
È una densità: è $\geq 0$, e l'area sotto il triangolo (base 2, altezza 1) vale $\frac{2 \cdot 1}{2} = 1$. Calcoliamo la
CDF integrando un pezzo alla volta.
\begin{itemize}
  \item Per $x < -1$ a sinistra non c'è area: $F_X(x) = 0$.
  \item Per $-1 \leq x \leq 0$ la densità vale $1 + t$, con primitiva $t + \frac{t^2}{2}$:
        \[ F_X(x) = \int_{-1}^{x} (1 + t) \dd t = \left[ t + \tfrac{t^2}{2} \right]_{-1}^{x} = x + \tfrac{x^2}{2} -
           \left( -1 + \tfrac12 \right) = \frac{(1 + x)^2}{2} . \]
  \item Per $0 \leq x \leq 1$ all'area di sinistra, $F_X(0) = \frac12$, si aggiunge quella fra 0 e $x$, dove la densità
        vale $1 - t$: $\int_0^x (1 - t) \dd t = x - \frac{x^2}{2}$. Quindi $F_X(x) = \frac12 + x - \frac{x^2}{2} =
        1 - \frac{(1 - x)^2}{2}$.
  \item Per $x > 1$ tutta l'area è a sinistra: $F_X(x) = 1$.
\end{itemize}
\begin{center}
\begin{tikzpicture}
\begin{groupplot}[group style={group size=2 by 1, horizontal sep=1.6cm},
  quaderno, width=6.6cm, height=4.4cm, xmin=-1.8, xmax=1.9, ymin=0, ymax=1.25, xtick={-1,-0.5,0.5,1},
  ytick={0.125,0.5,0.875,1}, yticklabels={$\frac18$,$\frac12$,$\frac78$,$1$}]
\nextgroupplot[ylabel=$f_X(x)$, xlabel=$x$]
  \addplot[fill=evid, draw=none, domain=-0.5:0.5, samples=3] {1-abs(x)} \closedcycle;
  \addplot[mspiA, domain=-1.6:1.6, samples=161] {max(0,1-abs(x))};
\nextgroupplot[ylabel=$F_X(x)$, xlabel=$x$]
  \addplot[mspiB, domain=-1.6:1.6, samples=161] {x<-1 ? 0 : (x<0 ? (1+x)^2/2 : (x<1 ? 1-(1-x)^2/2 : 1))};
  \draw[mspigray, dashed] (axis cs:-0.5,0) -- (axis cs:-0.5,0.125) -- (axis cs:-1.8,0.125);
  \draw[mspigray, dashed] (axis cs:0.5,0) -- (axis cs:0.5,0.875) -- (axis cs:-1.8,0.875);
\end{groupplot}
\end{tikzpicture}
\end{center}
La probabilità che $X$ cada fra $-\frac12$ e $\frac12$ è l'area gialla a sinistra, oppure la differenza dei valori della
CDF a destra:
\[ P\left(-\tfrac12 \leq X \leq \tfrac12\right) = F_X\left(\tfrac12\right) - F_X\left(-\tfrac12\right) = \left( 1 -
   \tfrac18 \right) - \tfrac18 = \tfrac34 . \]
Verifica geometrica: fuori dalla zona gialla restano due triangolini con base e altezza $\frac12$, di area $\frac18$
ciascuno, e $1 - \frac18 - \frac18 = \frac34$. Si noti anche che la CDF cresce più in fretta dove la densità è più alta
(attorno a 0): la sua pendenza è proprio la densità.
\end{esempio}

\subsection{La CDF delle variabili discrete}

La CDF si può definire anche per le variabili discrete (e per quelle «miste», di cui non parliamo): $F_X(x) =
P(X \leq x)$ ha senso per qualsiasi variabile aleatoria. In questo caso, però, la proprietà di continuità va
«rimodulata»: la CDF di una variabile discreta è una \textbf{scala}, piatta fra un valore dell'alfabeto e il successivo,
con un \textbf{salto} in ogni valore $x$ alto esattamente $p_X(x)$.

\begin{figure}[H]
\centering
\begin{tikzpicture}
\begin{axis}[quaderno, width=10cm, height=4.6cm, xmin=-0.5, xmax=7.6, ymin=0, ymax=1.18, xtick={1,...,6},
  ytick={0.1667,0.3333,0.5,0.6667,0.8333,1}, yticklabels={$\frac16$,$\frac26$,$\frac36$,$\frac46$,$\frac56$,$1$},
  xlabel=$x$, ylabel=$F_X(x)$]
  \addplot[mspiA, line width=1.3pt, jump mark left] coordinates {(-0.5,0) (1,0.1667) (2,0.3333) (3,0.5) (4,0.6667) (5,0.8333) (6,1) (7.4,1)};
  \addplot[only marks, mark=*, mark size=1.8pt, mspiA] coordinates {(1,0.1667) (2,0.3333) (3,0.5) (4,0.6667) (5,0.8333) (6,1)};
  \addplot[only marks, mark=o, mark size=1.8pt, mspiA, fill=white] coordinates {(1,0) (2,0.1667) (3,0.3333) (4,0.5) (5,0.6667) (6,0.8333)};
\end{axis}
\end{tikzpicture}
\caption{La CDF del lancio di un dado: in ogni faccia sale di un gradino alto $\frac16$. Il pallino pieno indica il
valore nel punto di salto: $F_X(3) = P(X \leq 3) = \frac36$, perché il 3 è compreso.}
\end{figure}

% ============================================================
\section{La media di una variabile continua}
% ============================================================

\dfn{Media statistica di una variabile continua}{
  La \term{media statistica} di una variabile aleatoria continua con densità $f_X(x)$ è
  \[ E[X] = \mu_X = \int_{-\infty}^{\infty} x\, f_X(x) \dd x , \]
  purché l'integrale converga.
}

La definizione è la stessa del caso discreto, $\sum_x x\, p_X(x)$, con l'integrale al posto della somma e
$f_X(x) \dd x$ al posto di $p_X(x)$. Per giustificarla in modo convincente si può approssimare la variabile continua
con una discreta.

\begin{dimostrazione}[la media come limite di una variabile quantizzata]
\passo{Si quantizza $X$} Si divide l'asse reale in intervalli di ampiezza $\Delta$ e si definisce la versione
«quantizzata» di $X$: se $X$ cade nell'intervallo $[i\Delta, (i+1)\Delta)$, la si sostituisce con un valore $x_i$ di
quell'intervallo,
\[ X^{\Delta} = x_i \in [i\Delta, (i+1)\Delta) \qquad \text{se} \quad i\Delta \leq X < (i+1)\Delta . \]
È una variabile discreta, con probabilità
\[ P(X^{\Delta} = x_i) = \int_{i\Delta}^{(i+1)\Delta} f_X(x) \dd x . \]
\passo{La media della variabile discreta} Con la definizione del Capitolo~5:
\[ E[X^{\Delta}] = \sum_{i=-\infty}^{\infty} x_i\, P(X^{\Delta} = x_i) = \sum_{i=-\infty}^{\infty} x_i
   \int_{i\Delta}^{(i+1)\Delta} f_X(x) \dd x . \]
\passo{Si stringono gli intervalli} Quando $\Delta$ è piccolo, l'area sotto $f_X$ su un intervallo è circa
l'altezza per la base, $f_X(x_i)\,\Delta$, e la quantizzazione diventa sempre più fine. Se $x f_X(x)$ è integrabile
(secondo Riemann), la somma diventa una somma di Riemann della funzione $x f_X(x)$:
\[ E[X] = \lim_{\Delta \to 0} E[X^{\Delta}] = \lim_{\Delta \to 0} \sum_{i=-\infty}^{\infty} x_i\, f_X(x_i)\,\Delta =
   \int_{-\infty}^{\infty} x\, f_X(x) \dd x . \]
\end{dimostrazione}

\begin{figure}[H]
\centering
\begin{tikzpicture}
\begin{axis}[quaderno, width=11cm, height=4.6cm, xmin=-1.3, xmax=1.4, ymin=0, ymax=1.15, xtick={-1,-0.5,0,0.5,1},
  ytick=\empty, xlabel=$x$, grid=none]
  \addplot[ybar interval, fill=mspiA!18, draw=mspiA!50, samples=9, domain=-1:1] {1-abs(x+0.125)};
  \addplot[mspiA, domain=-1.2:1.2, samples=121] {max(0,1-abs(x))};
  \addplot[only marks, mark=*, mark size=1.8pt, mspiB] coordinates {(-0.875,0) (-0.625,0) (-0.375,0) (-0.125,0) (0.125,0) (0.375,0) (0.625,0) (0.875,0)};
  \node[font=\scriptsize, mspiB!85!black, anchor=south west, align=left] at (axis cs:0.95,0.45)
    {tutta la probabilità di un\\ intervallo viene spostata\\ nel suo punto $x_i$};
\end{axis}
\end{tikzpicture}
\caption{La quantizzazione della densità triangolare con $\Delta = \frac14$: ogni rettangolo azzurro ha area circa
uguale alla probabilità del suo intervallo, che la variabile quantizzata $X^{\Delta}$ concentra nel punto arancione.}
\end{figure}

\begin{esempio}{Due medie}
\begin{itemize}
  \item La densità triangolare è simmetrica attorno a 0, e la sua media è 0: i contributi $x f_X(x)$ a destra
        (positivi) e a sinistra (negativi) sono uguali e opposti.
  \item Sia $f_X(x) = 2x$ per $0 \leq x \leq 1$ (e 0 altrove): una densità che cresce, quindi i valori grandi sono più
        probabili. È una densità perché $\int_0^1 2x \dd x = \left[ x^2 \right]_0^1 = 1$. La media è
        \[ E[X] = \int_0^1 x \cdot 2x \dd x = \int_0^1 2x^2 \dd x = \left[ \tfrac{2x^3}{3} \right]_0^1 = \tfrac23 ,
        \]
        spostata verso destra rispetto al punto di mezzo $\frac12$, come ci si aspetta.
\end{itemize}
Anche nel continuo la media è il \textbf{baricentro}: una sbarra con densità di massa $f_X(x)$ starebbe in equilibrio
appoggiata nel punto $E[X]$.
\end{esempio}

\gbox{In sintesi}{mspiA!80!black}{
\begin{itemize}
  \item Per una variabile continua ogni singolo valore ha probabilità zero: gli intervalli, non i punti, hanno
        probabilità. Solo un'infinità al più numerabile di valori può avere probabilità positiva.
  \item La \textbf{densità} $f_X(x) = \lim P(x - \frac{\Delta x}{2} \leq X \leq x + \frac{\Delta x}{2})/\Delta x$ è una
        probabilità per unità di lunghezza: $f_X \geq 0$, $\int f_X = 1$, e può superare 1.
  \item $P(a_1 \leq X \leq a_2) = \int_{a_1}^{a_2} f_X$: l'\textbf{integrale} è l'area sotto la curva, limite delle somme
        di Riemann; si calcola con una primitiva, $\int_a^b f = G(b) - G(a)$.
  \item La \textbf{CDF} $F_X(x) = P(X \leq x) = \int_{-\infty}^{x} f_X$ è continua, crescente, va da 0 a 1, e
        $f_X = F_X'$ (teorema fondamentale del calcolo integrale). La CCDF è $1 - F_X$. Per le variabili discrete la CDF
        è una scala.
  \item \textbf{Media}: $E[X] = \int x\, f_X(x) \dd x$, limite della media di una versione quantizzata di $X$.
\end{itemize}
}
